\chapter{Contexte, problématique et état de l'art}

\section*{Introduction}
Ce premier chapitre pose les fondations du projet. Nous y présentons le contexte général dans lequel s'inscrit GREENLY, la problématique qui a motivé sa conception, les objectifs poursuivis, ainsi qu'une étude de l'état de l'art des solutions existantes de calcul d'empreinte carbone. Une comparaison critique de ces solutions permet enfin de justifier les choix qui ont guidé la conception de notre propre plateforme.

\section{Contexte}

Selon le Groupe d'experts intergouvernemental sur l'évolution du climat (GIEC), la limitation du réchauffement climatique à 1,5°C exige une réduction drastique et rapide des émissions mondiales de gaz à effet de serre. Si une part significative de ces émissions provient de l'industrie lourde et des transports collectifs, une proportion tout aussi importante résulte des choix individuels quotidiens~: déplacements motorisés, consommation énergétique domestique, habitudes alimentaires, gestion des déchets ménagers et usage des équipements électroniques.

Prendre conscience de son propre impact environnemental est une première étape indispensable pour amorcer un changement de comportement. C'est le rôle des \textit{calculateurs d'empreinte carbone}~: des outils, le plus souvent numériques, permettant à un individu de renseigner ses habitudes de consommation afin d'obtenir une estimation chiffrée de son impact en équivalent CO2 (CO2e).

Le développement rapide des techniques de Machine Learning et la disponibilité croissante de jeux de données publics ouverts (\textit{open data}) — publiés par des organismes tels que l'\textit{Environmental Protection Agency} (EPA) américaine, l'université UCI ou la base collaborative \textit{Our World in Data} — offrent aujourd'hui une opportunité nouvelle~: dépasser les calculateurs à formules fixes pour construire des systèmes prédictifs capables d'apprendre des relations complexes entre les habitudes de consommation et leur impact réel, tout en offrant un retour visuel et analytique digne des outils de Business Intelligence professionnels.

\section{Problématique}

L'analyse des outils existants (détaillée en section~\ref{sec:etat-art}) révèle plusieurs limites récurrentes~:

\begin{itemize}[label=\textbullet,font=\normalsize]
    \item \textbf{Rigidité du modèle de calcul}~: la quasi-totalité des calculateurs grand public appliquent des formules linéaires fixes (facteur d'émission $\times$ quantité), incapables de capter les interactions non linéaires entre plusieurs variables (par exemple l'effet combiné du type de carburant, du type de véhicule et de la distance parcourue).
    \item \textbf{Absence de suivi dans le temps}~: la plupart des outils proposent un calcul ponctuel, sans historique ni tableau de bord permettant de suivre l'évolution de l'empreinte d'un utilisateur.
    \item \textbf{Absence de prévision}~: aucun outil grand public ne propose de projection statistique de la tendance future d'un utilisateur à partir de son historique.
    \item \textbf{Recommandations génériques}~: les conseils fournis sont, en général, des messages statiques identiques pour tous les utilisateurs, sans lien avec les facteurs qui pèsent réellement le plus dans le calcul de l'utilisateur concerné.
    \item \textbf{Manque de transparence méthodologique}~: les facteurs d'émission utilisés sont rarement documentés ni reliés à des sources vérifiables.
\end{itemize}

La problématique centrale de ce projet peut donc se formuler ainsi~: \emph{comment concevoir une plateforme de calcul d'empreinte environnementale qui soit à la fois précise (fondée sur de véritables modèles prédictifs entraînés sur des données réelles), transparente (facteurs d'émission documentés et traçables), personnalisée (recommandations fondées sur l'importance réelle des variables) et outillée pour le suivi dans la durée (tableaux de bord et prévisions~)?}

\section{Objectifs}

Le projet GREENLY poursuit les objectifs suivants~:

\begin{enumerate}
    \item Concevoir et entraîner des modèles de Machine Learning capables de prédire l'empreinte environnementale (carbone ou hydrique) pour six domaines d'usage~: transport, énergie domestique, électronique, eau, alimentation et déchets.
    \item Fonder ces modèles sur des données publiques réelles et documentées, plutôt que sur des constantes arbitraires, afin de garantir la traçabilité scientifique des résultats.
    \item Comparer objectivement plusieurs algorithmes d'apprentissage supervisé (Random Forest, Gradient Boosting, XGBoost, LightGBM, CatBoost) et sélectionner automatiquement le plus performant pour chaque domaine.
    \item Développer une interface web moderne, ergonomique et réactive permettant à l'utilisateur de saisir ses données et de visualiser instantanément ses résultats.
    \item Intégrer un module de Business Intelligence~: indicateurs clés de performance, tableaux de bord interactifs et visualisations multi-dimensionnelles de l'historique utilisateur.
    \item Développer un module de prévision statistique permettant de projeter la tendance future de l'empreinte d'un utilisateur à partir de son historique réel.
    \item Générer des recommandations personnalisées fondées sur l'importance réelle des variables (\textit{feature importance}) telle que calculée par les modèles entraînés, plutôt que sur des messages génériques.
    \item Garantir la sécurité et la confidentialité des données utilisateurs via une architecture d'authentification et de contrôle d'accès robuste.
\end{enumerate}

\section{État de l'art}
\label{sec:etat-art}

Plusieurs outils de calcul d'empreinte carbone sont aujourd'hui accessibles au grand public. Cette section en présente les principaux représentants.

\subsection{Nos Gestes Climat (ADEME)}
Développé par l'Agence de la transition écologique (ADEME) française, \textit{Nos Gestes Climat} \cite{nosGestesClimat} est l'un des calculateurs les plus utilisés en France. Il repose sur un long questionnaire couvrant le logement, le transport, l'alimentation, la consommation et les services publics, et applique des facteurs d'émission officiels (Base Carbone\textsuperscript{\textregistered}) pour produire une estimation annuelle en tonnes de CO2e.

\subsection{WWF Footprint Calculator}
Le calculateur du \textit{World Wide Fund for Nature} \cite{wwfFootprint} propose une estimation rapide de l'empreinte écologique globale (et non uniquement carbone) à partir d'un questionnaire simplifié, avec des résultats exprimés en nombre de « planètes » nécessaires si tout le monde vivait comme l'utilisateur.

\subsection{EPA Household Carbon Footprint Calculator}
Proposé par l'agence américaine de protection de l'environnement \cite{epaCalculator}, cet outil applique directement les facteurs d'émission officiels EPA (eGRID pour l'électricité, facteurs de combustion pour le gaz et le fioul) à des données saisies par l'utilisateur (factures d'électricité, kilométrage annuel, etc.).

\subsection{Carbon Trust Footprint Calculator}
Le \textit{Carbon Trust} \cite{carbonTrust}, organisme britannique spécialisé dans l'accompagnement bas-carbone, propose un calculateur orienté ménages et petites entreprises, avec un focus sur les recommandations de réduction associées à chaque poste d'émission.

\subsection{Approches académiques par apprentissage automatique}
Au-delà des outils grand public, plusieurs travaux de recherche explorent l'usage du Machine Learning pour l'estimation d'empreinte carbone dans des contextes spécifiques~: prédiction des émissions liées au trafic routier à partir de modèles d'ensemble d'arbres de décision \cite{scikitlearn}, ou encore modélisation de la consommation énergétique résidentielle à partir de données de capteurs, comme dans le jeu de données \textit{Appliances Energy Prediction} de l'UCI Machine Learning Repository \cite{candanedo2017}, utilisé comme fondation réelle du module « Énergie » de GREENLY (voir chapitre~3). Ces travaux démontrent la capacité des modèles à base d'arbres (Random Forest, Gradient Boosting) à capter des relations non linéaires que les formules fixes ne peuvent représenter, mais restent, pour la plupart, cantonnés à un seul domaine d'application et ne débouchent pas sur une plateforme utilisateur complète intégrant suivi et prévision.

\section{Comparaison des solutions}

Le tableau~\ref{tab:comparaison-solutions} synthétise une comparaison entre les principaux outils existants et la solution proposée, selon des critères clés identifiés à partir de la problématique.

\begin{table}[H]
\centering
\caption{Comparaison des solutions existantes de calcul d'empreinte carbone}
\label{tab:comparaison-solutions}
\begin{tabular}{|p{3.1cm}|c|c|c|c|c|}
\hline
\rowcolor{isiBlue!15}
\textbf{Critère} & \textbf{Nos Gestes Climat} & \textbf{WWF} & \textbf{EPA} & \textbf{Carbon Trust} & \textbf{GREENLY} \\
\hline
Approche de calcul & Formule fixe & Formule fixe & Formule fixe & Formule fixe & \textbf{Machine Learning} \\
\hline
Facteurs d'émission documentés & Oui & Partiel & Oui & Partiel & \textbf{Oui (sourcés)} \\
\hline
Multi-domaine (6 catégories) & Oui & Oui & Non (énergie/transport) & Oui & \textbf{Oui} \\
\hline
Historique \& suivi temporel & Non & Non & Non & Non & \textbf{Oui} \\
\hline
Tableau de bord BI & Non & Non & Non & Non & \textbf{Oui} \\
\hline
Prévision statistique & Non & Non & Non & Non & \textbf{Oui} \\
\hline
Score comparatif (percentile) & Non & Non & Non & Non & \textbf{Oui (Green Score)} \\
\hline
Recommandations personnalisées (feature importance) & Non & Non & Non & Générique & \textbf{Oui (dérivées du modèle)} \\
\hline
Ouvert / auto-hébergeable & Oui (open-source) & Non & Non & Non & \textbf{Oui} \\
\hline
\end{tabular}
\end{table}

Cette comparaison met en évidence le positionnement singulier de GREENLY~: c'est, parmi les solutions étudiées, la seule à combiner une véritable approche prédictive par apprentissage automatique, une architecture de suivi temporel outillée par un tableau de bord de Business Intelligence, et un module de prévision statistique des tendances futures de l'utilisateur.

\section*{Conclusion}
Ce chapitre a permis de situer le projet GREENLY dans son contexte, de formuler la problématique à laquelle il répond et d'en dégager les objectifs. L'étude de l'état de l'art a montré que si plusieurs calculateurs d'empreinte carbone existent déjà, aucun ne combine approche prédictive par Machine Learning, suivi historique et prévision statistique au sein d'une même plateforme. Le chapitre suivant détaille la spécification des besoins fonctionnels et non fonctionnels du système, ainsi que la démarche méthodologique Agile Scrum adoptée pour la conduite du projet.
